% Folder 70, batch 1 (archivists' pages 1-20).
% Pass: Opus 5.5 (claude-opus-5-5), 2026-09-29 — first pass, unchecked against the pages by a human.
%
% Pages 1 and 2 are covers inscribed in pencil; their words become the section
% headings (no \page). Pages 3-20 are a single run on regular n-polyhedra and
% their affine realisations. Notation fixed for the batch: \mathfrak{G}_n for
% his script G (the n-cartographic group), \underline{R} for the underlined R
% (the homogeneous space of "repères"), f_i for the flag's i-facets, s_i for
% the vertices, e_i = s_i - s_0, \lambda_{ij} for the parameters.
\documentclass[11pt,a4paper]{article}
\input{../preamble/grothendieck}

\folder{70}
\batch{1}
\pages{1}{20}
\foldertitle{Réalisations géométriques de structures combinatoires (n-polyèdres, n-hyperpolyèdres [2-polyèdres réguliers]…) (Vieilles rédactions) : notes manuscrites (s.d.).}
\dating{[à partir de 1976]}
\watermark{Édition de démonstration}

\begin{document}

\section{$n$-polyèdres et $n$-hyperpolyèdres}
\note{Titre repris de la couverture (p.~1 des archivistes), au crayon, de sa main selon toute apparence ; le feuillet ne porte rien d'autre.}

\section{$n$-polyèdres réguliers en coordonnées affines universelles (vieilles rédactions)}
\note{Titre repris de la seconde couverture (p.~2), au crayon ; elle porte aussi la mention « (23~p) ». Les pp.~3--20 de ce lot ne font que dix-huit feuillets : la liasse annoncée se poursuit au-delà du lot.}

\page{3}
\emph{$n$-polyèdres réguliers.}

Considérons \struck{\ill{}} une $n$-carte \emph{régulière}, définie par un espace homogène $\underline{R}$ \add{(repères)} \struck{\uncertain{sous}} sous le groupe $n$-cartographique
\[
\mathfrak{G}_n = \Bigl\{ \sigma_0,\dots,\sigma_n \Bigm| \sigma_0^2=\dots=\sigma_n^2=1,\quad \sigma_i\sigma_j=\sigma_j\sigma_i \ \text{si}\ \begin{cases} j\geq i+2 \\ 0\leq i,j\leq n \end{cases} \Bigr\}
\]
\note{la notation $\mathfrak{G}_n$ rend un G cursif bouclé, le même d'un bout à l'autre du lot.}

tel que le stabilisateur \add{$\mathfrak{G}_{n,\underline{r}}$} \struck{de} \uncertain{d'un} $\underline{r}\in\underline{R}$ soit distingué dans $\mathfrak{G}_n$, donc indépendant de $\underline{r}$, soit $H$, \add{et $\mathfrak{G}_n/H=\Gamma$}. On suppose que pour toute partie $I$ de $[0,n]\cap\mathbb{N}$ ne contenant pas deux entiers consécutifs, le produit \struck{de} $\sigma_I=\prod_{i\in I}\sigma_i\notin H$, i.e. le produit $\sigma_H$ \note{sic : $\sigma_H$ pour $\sigma_I$} des $\sigma_i$ \emph{dans} $\Gamma$ est $\neq 1$.
\marginal{\uncertain{inutile} \ldots}

On pose
\[
\begin{cases}
p_1 = \text{ordre de } \sigma_0\sigma_1 \\
\dots \\
p_n = \text{ordre de } \sigma_{n-1}\sigma_n
\end{cases}
\qquad \text{\emph{NB} les $p_i$ peuvent être infinis !}
\]
et on suppose
\[
3\leq p_i\leq +\infty \qquad \text{pour } 1\leq i\leq n .
\]
On suppose ici que $\Gamma$ est le groupe engendré par les $\sigma_i$, avec (en plus des relations plus haut définissant $\mathfrak{G}_n$) les relations supplémentaires
\[
(\sigma_{i-1}\sigma_i)^{p_i}=1 .
\]

\page{4}
Si
\[
p_* = (p_1,\dots,p_n)\in\bigl([3,+\infty]\cap\mathbb{N}\bigr)^n ,
\]
\struck{\ill{}} $\Gamma=\Gamma_{p_*}$ est entièrement déterminé par \struck{$p_*$} $p_*$ et, quel que soit $p_*$, les conditions \struck{\ill{}}
\[
\sigma_I\neq 1 \ \text{dans}\ \Gamma_{p_*}
\]
sont satisfaites dans $\Gamma=\Gamma_{p_*}$. Tous les $\underline{R}$ sont isom. entre eux, et quand on prend des $\underline{R}$ munis d'un « repère » $\underline{r}\in\underline{R}$, ils sont canoniquement isom. entre eux, et \add{com.} isom. à
\[
\underline{R}_{p_*} = \mathfrak{G}_n/\Gamma_{p_*}
\]
\note{écrit ainsi ; d'après la p.~3 on attendrait $\mathfrak{G}_n/H$, c'est-à-dire $\Gamma_{p_*}$ lui-même comme ensemble.}
muni du repère image de $1\in\mathfrak{G}_n$.

Pour toute partie $H$ de \struck{$\mathfrak{G}$} $[0,n]\cap\mathbb{N}$, on désigne par $\mathfrak{G}_{n,H}$ le sous-groupe de $\mathfrak{G}_n$ engendré par les $\sigma_i$, $i\notin H$. Donc
\[
H\subset K \Longrightarrow \mathfrak{G}_K\subset\mathfrak{G}_H ,
\qquad
\underbrace{\underline{R}/\mathfrak{G}_K}_{=D_K} \longrightarrow \underbrace{\underline{R}/\mathfrak{G}_H}_{=D_H} .
\]
Si $H=\{i\}$, on trouve
\[
D_{\{i\}} = D_i = \text{ens. des $i$-facettes \add{du polyèdre régulier $\Pi$ défini par $\underline{R}$}}.
\]

\page{5}
Sauf erreur, on est dans les conditions qu'il faut pour que la \struck{polyèdre comb.} \add{carte comb.} régulière définie par le \uncertain{singulet} $\underline{R}\simeq\underline{R}_{p_*}$ soit définie par une \emph{géométrie d'incidence} ayant \struck{\ill{}} $D_i$ comme ens. des facettes \add{$\Pi$} de dim.~$i$ (donc soit un polyèdre comb. \emph{régulier}).

Une \emph{réalisation géom.} \add{affine régulière} (de $\Pi$ \struck{\ill{}} sur un \struck{\ill{}} anneau \ill{} comme le \ill{}) est la donnée
\begin{itemize}
  \item[a)] d'un torseur (espace affine) \add{$E$} sous un $k$-Module loc. libre de rang $n+1$ ;
  \item[b)] d'applications
  \[
  \varphi_i : D_i \longrightarrow \text{ens. des sous-espaces affines de dim.~$i$ dans $E$} \qquad (0\leq i\leq n)
  \]
  respectant les relations d'incidence,
  \item[\struck{c)}] \struck{d'un homomorphisme de groupes} avec condition supplémentaire \struck{\ill{}} :
  \struck{pour tout \ill{} $\mathrm{Aut}(\Pi)\to\mathrm{Aut\,aff}(E)$ compatible avec les}
  \[
  \forall g\in\mathrm{Aut}(\Pi),\ \exists!\ \bar g\in\mathrm{Aff}(E)\ \text{tel que}\quad \bar g\circ\varphi_i=\varphi_i\circ g_{D_i}\quad \forall\, 0\leq i\leq n
  \]
  --- cette unicité reste vraie fibre par fibre\ldots
\end{itemize}

\page{6}
\marginal{: vérifier}
Il revient au \uncertain{même} de se donner une application
\[
\varphi_{\underline{R}} : \underline{R} \longrightarrow \struck{\mathrm{Rep\,aff}(E)}\ \mathrm{Drap.aff}_n(E)
\]
\add{(drapeaux maximaux de $E$)} \struck{et un homo \ill{} qui \ill{} un unique} et un homo
\[
\mathfrak{G} \xrightarrow{\ \varphi_{\mathfrak{G}}\ } \mathrm{Aut\,aff}\,E
\]
compatible avec $\varphi_{\underline{R}}$, tel que les sommets des repères drapeaux de l'image de $\underline{R}$ engendrent affinement $E$. (Ce qui implique l'unicité de $\varphi_{\mathfrak{G}}$ \add{éliminons} \ill{} éliminant\ldots), et la condition suivante : \struck{\ill{} le sous-groupe $\mathfrak{G}_{\underline{r}}$ de $\mathfrak{G}$ stabilisateur de la composante $\underline{r}\in\underline{R}$, opérant sur \ill{} comme \ill{}} \ill{} $\mathfrak{G}$, il suffit en fait, étant choisi \add{grâce à l'identification $\underline{R}\simeq\underline{R}_{p_*}$} une origine $\underline{r}\in\underline{R}$, de se donner $\varphi_{\mathfrak{G}}$ et $\varphi_{\underline{R}}(\underline{r})=r\in\mathrm{Drap.aff}_n(E)$, avec comme seules conditions que
\begin{itemize}
  \item[a)] si $j\neq i$, $\sigma_j$ invarie \uncertain{$f_i$} ;
  \item[b)] les $g\cdot f_0$ ($g\in\mathfrak{G}=\Gamma$) engendrent affinement $E$ fibre par fibre\ldots
\end{itemize}
\marginal{\uncertain{$E$}, \ill{} \ill{} \ill{} de $\varphi$}
\note{la marge gauche porte, en travers, une addition de quatre lignes que nous ne lisons pas ; « $\varphi_{\mathfrak{G}} = \rho$ », sous $\varphi_{\mathfrak{G}}$, est une lecture incertaine d'un ajout interlinéaire.}

\page{7}
Posons
\[
\begin{cases}
s_0 = f_0 \\
s_1 = \sigma_0 f_0 \\
s_2 = \sigma_1 s_1 \\
\dots \\
s_i = \sigma_{i-1} s_{i-1} \\
\dots \\
s_{n+1} = \sigma_n s_n
\end{cases}
\]
Je dis que les $(s_i)_{0\leq i\leq n+1}$ forment une base affine de $E$ fibre par fibre, \add{et que $f_i=$ env. aff. de $(s_0,\dots,s_i)$, $0\leq i\leq n$} \struck{\ill{}}. Supposons prouvé, pour $i$ donné, \add{$0\leq i\leq n$,} que $\{s_0,\dots,s_i\}$ \add{est} affinement libre fibre par fibre, \add{que $f_i=$ env. aff. de $\{s_0,\dots,s_i\}$}, \struck{$0\leq i\leq n+1$ \ill{}} et que $\sigma_k s_j=s_j$ si \struck{$j\leq i$} $k>j$ \add{(pour $j\leq i$)}. Prouvons que $s_{i+1}=\sigma_i s_i$ n'est pas dans \add{l'env. aff.} $f_i$ de $\{s_0,\dots,s_i\}$, et que $f_{i+1}$ est l'env. aff. de $\{s_0,\dots,s_i,\sigma_i s_i\}$, puis que $\sigma_k s_{i+1}=s_{i+1}$ si $k>i+1$.

D'abord, comme $s_i\in f_i\subset f_{i+1}$ et $\sigma_i$ fixe $f_{i+1}$, on a $s_{i+1}=\sigma_i(s_i)\in f_{i+1}$, donc il faut seulement prouver \add{(\ill{} \ill{} sur les fibres)} que $\sigma_i(s_i)\notin f_i$. \struck{Mais si} Mais si on avait $\sigma_i(s_i)\in f_i$, \struck{comme $\sigma_i f_i\ldots f_{i+1}$, donc $\sigma_i$ \ill{} on aurait $\sigma_k f_i=f_i$ pour $k\neq i$,} \struck{on aurait} on aurait
\[
\sigma_i(f_i) = \sigma_i\,\mathrm{Env}(f_{i-1},s_i) = \mathrm{Env}\bigl(\underbrace{\sigma_i f_{i-1}}_{f_{i-1}},\ \underbrace{\sigma_i s_i}_{\in f_i}\bigr)\subset f_i
\]
\note{la fin de la première ligne de la formule, raturée, n'est pas lue : \struck{\ill{}}.}
donc $\sigma_i(f_i)=f_i$, donc comme $\sigma_j f_i=f_i$ pour $j\neq i$, on aurait $\sigma_k f_i=f_i$ pour $0\leq k\leq n$, donc $f_i$ fixé par $\Gamma=\mathfrak{G}$, donc

\page{8}
(comme $s_0\in f_i$) $\mathfrak{G}\cdot s_0\subset f_i$, contrairement à l'hyp. que $\mathfrak{G}\cdot s_0$ engendre l'espace fibre par fibre. Il reste à prouver que
\[
\sigma_k s_{i+1}=s_{i+1}\quad\text{pour } n\geq k>i+1,\ \text{i.e.}\ \sigma_k\sigma_i s_i=\sigma_i s_i ,
\]
or
\[
\sigma_k\sigma_i s_i = \sigma_i\underbrace{\sigma_k s_i}_{=s_i}\ (\text{car } k>i+1) = \sigma_i s_i ,\quad \text{OK.}
\]
\note{un double trait souligne ce qui précède ; au-dessous, un mot \ill{}.}

\[
\begin{cases}
\{s_i\}_{0\leq i\leq n+1}\ \text{base affine fibre par fibre} \\
f_i=\mathrm{Env}(s_0,\dots,s_i) \quad \text{si } 0\leq i\leq n+1 \quad (\text{posant } f_{n+1}=E) \\
\sigma_k s_i = s_i \quad \text{si } \struck{k}\ 0\leq i<k\leq n \\
\sigma_i s_i = s_{i+1} \quad 0\leq i\leq n
\end{cases}
\]
impliquent déjà \struck{que} formellement que
\[
\sigma_j f_i = f_i \quad\text{si } j>i ,
\]
donc il faut encore voir \uncertain{où} \ill{} que \struck{(clairement si $j>i$, et si $j<i$)}
\[
\sigma_j f_i = f_i \quad\text{si } j<i .
\]
\struck{Sont} Plus les relations (\ill{} \ill{})
\[
\begin{cases}
\sigma_i^2=1 \quad \struck{\ill{}} & 0\leq i\leq n \\
\sigma_i\sigma_j=\sigma_j\sigma_i & \text{si } 0\leq i\ \ill{} \\
(\sigma_{i-1}\sigma_i)^{p_i}=1 & 1\leq i\ \ill{}
\end{cases}
\]

\page{9}
\[
e_i = s_i - s_0 \in V, \quad 1\leq i\leq n+1 \ \text{base de } V, \qquad s_0 \ \text{\emph{origine} de } E
\]
\[
\begin{cases}
\sigma_0 s_0 = s_1 \\
\sigma_0 s_1 = s_0 \\
\sigma_0 s_2 = s_2^0 \in f_2\smallsetminus f_1 \\
\sigma_0 s_3 = s_3^0 \in f_3\smallsetminus f_2 \\
\dots \\
\sigma_0 s_{n+1} = s_{n+1}^0 \in f_{n+1}\smallsetminus f_n
\end{cases}
\qquad
\begin{aligned}
s_2^0 &= s_0+\lambda_2^{01}e_1+\lambda_2^{02}e_2 \\
s_3^0 &= s_0+\lambda_3^{01}e_1+\lambda_3^{02}e_2+\lambda_3^{0\,\uncertain{3}}e_3 \\
&\dots \\
s_{n+1}^0 &= s_0+\lambda_{n+1}^{01}e_1+\lambda_{n+1}^{02}e_2+\dots+\lambda_{n+1}^{0,n+1}e_{n+1}
\end{aligned}
\]
\[
\begin{cases}
\sigma_1 s_0 = s_0 \\
\sigma_1 s_1 = s_2 \\
\sigma_1 s_2 = s_1 \\
\sigma_1 s_3 = s_3^1 \in f_3\smallsetminus f_2 \\
\sigma_1 s_4 = s_4^1 \in f_4\smallsetminus f_2 \\
\dots \\
\sigma_1 s_{n+1} = s_{n+1}^1 \in f_{n+1}\smallsetminus f_n
\end{cases}
\qquad
\begin{cases}
\sigma_2 s_0 = s_0 \\
\sigma_2 s_1 = s_1 \\
\sigma_2 s_2 = s_3 \\
\sigma_2 s_3 = s_2 \\
\sigma_2 s_4 = s_4^2 \in f_4\smallsetminus f_3 \\
\dots \\
\sigma_2 s_{n+1} = s_{n+1}^2 \in f_{n+1}\smallsetminus f_n
\end{cases}
\]
\note{« $f_4\smallsetminus f_2$ » pour $\sigma_1 s_4$ est écrit ainsi ; le schéma des autres lignes donnerait $f_4\smallsetminus f_3$.}
\[
\begin{cases}
\sigma_i s_0 = s_0 \\
\sigma_i s_1 = s_1 \\
\dots \\
\sigma_i s_{i-1} = s_{i-1} \\
\sigma_i s_i = s_{i+1} \\
\sigma_i s_{i+1} = s_i \\
\sigma_i s_{i+2} = s_{i+2}^i \in f_{i+2}\smallsetminus f_{i+1} \\
\dots \\
\sigma_i s_{n+1} = s_{n+1}^i \in f_{n+1}\smallsetminus f_n
\end{cases}
\quad 0\leq i\leq \struck{\ill{}}
\qquad
\begin{cases}
\sigma_n s_0 = s_0 \\
\sigma_n s_1 = s_1 \\
\dots \\
\sigma_n s_{n-1} = s_{n-1} \\
\sigma_n s_n = s_{n+1} \\
\sigma_n s_{n+1} = s_n
\end{cases}
\]

\page{10}
Compte tenu des relations de commutation, $\sigma_0,\dots,\sigma_n$ sont déterminés quand on se donne les \struck{\ill{}} valeurs critiques \struck{\ill{}}
\[
\begin{aligned}
\sigma_0 s_2 &= s_2^0\ (\in f_2) = s_0+\lambda_{01}e_1+\lambda_{02}e_2 \\
\sigma_1 s_3 &= s_3^1\ (\in f_3) = s_0+\lambda_{11}e_1+\lambda_{12}e_2+\lambda_{13}e_3 \\
&\dots \\
\sigma_{n-1} s_{n+1} &= s_{n+1}^{n-1}\ (\in f_{n+1}) = s_0+\lambda_{n-1,1}e_1+\lambda_{n-1,2}e_2+\dots+\lambda_{n-1,n+1}e_{n+1}
\end{aligned}
\]
défini par \uncertain{une} matrice à $n$ \emph{lignes} et $n+1$ colonnes
\[
\begin{pmatrix}
\lambda_{01} & \lambda_{02} & 0 & 0 & \cdots & 0 \\
\lambda_{11} & \lambda_{12} & \lambda_{13} & 0 & \cdots & 0 \\
\cdots & & & & & \\
\lambda_{n-1,1} & \lambda_{n-1,2} & \lambda_{n-1,3} & \lambda_{n-1,4} & \cdots & \lambda_{n-1,n+1}
\end{pmatrix}
\]
\note{une troisième ligne, commencée puis biffée, s'achevait sur $\lambda_{24}$.}
qui dépend \emph{a priori} de
\[
2+3+\dots+(n+1) = \frac{(n+2)(n+1)}{2}-1 = \frac{n(n+3)}{2}
\]
paramètres. En effet, on vérifie les
\[
\sigma_i s_j = \sigma_i\bigl(\sigma_{j-1}\,\struck{\sigma_{j-2}\cdots\sigma_{i+2}\,s_{i+2}}\bigr) = (\sigma_{j-1}\cdots\sigma_{i+2})\,\sigma_i s_{i+2}
\qquad j\geq i+3\ \text{i.e. } j-1\geq i+2
\]
\[
= (\sigma_{j-1}\sigma_{j-2}\cdots\sigma_{i+2})\,\underbrace{s_{i+2}^i}_{\in f_{i+2}} \qquad s_{i+2}^i = s_0+\lambda_{i1}e_1+\lambda_{i2}e_2+\dots+\lambda_{i,i+2}e_{i+2}
\]
\[
= s_0+\lambda_{i1}e_1+\dots+\lambda_{i,i+1}e_{i+1}+\lambda_{i,i+2}e_j
\]
(vrai aussi si $j=i+2$ donc)
\[
\boxed{\ \sigma_i s_j = (s_0+\lambda_{i1}e_1+\dots+\lambda_{i,i+1}e_{i+1})+\lambda_{i,i+2}e_j \quad\text{si } j\geq i+2\ }
\]

\page{11}
\[
\sigma_i e_j =
\begin{cases}
e_j & \text{si } j<i \\
e_{i+1} & \text{si } j=i \\
e_i & \text{si } j=i+1 \\
\lambda_{i1}e_1+\lambda_{i2}e_2+\dots+\lambda_{i,i+1}e_{i+1}+\lambda_{i,i+2}e_j & \text{si } j\geq i+2
\end{cases}
\qquad (i\geq 1)
\]
\[
\det(\sigma_i|_V) = -(\lambda_{i,i+2})^{n-i} =
\begin{cases}
-1 & \text{si } n-i \text{ pair} \quad \text{ok} \\
-\lambda_{i,i+2} & \text{si } n-i \text{ impair}
\end{cases}
\]
(\ill{} \add{donc} $\lambda_{i,i+2}=+1$)
\[
\sigma_0 e_j =
\begin{cases}
-e_1 & \text{si } j=1 \\
(\lambda_{01}-1)e_1+\lambda_{0,2}e_j & \text{si } j\geq 2
\end{cases}
\qquad \sigma_0 s_1
\]
\[
\sigma_0 e_j = \sigma_0 s_j-\sigma_0 s_0 = \sigma_0 s_j - s_1, \qquad \det(\sigma_0|_V) = -(\lambda_{0,2})^n
\]
\[
\sigma_0^2 s_j =
\begin{cases}
s_j & \text{si } j\leq 1 \\
\struck{\ill{}}\ \sigma_0(s_0+\lambda_{01}e_1+\lambda_{0,2}e_j)
\end{cases}
\]
\[
\struck{= s_0+e_1+}\qquad
s_0+e_1\ \uncertain{-}\ \lambda_{01}e_1+\lambda_{02}\bigl[(\lambda_{01}-1)e_1+\lambda_{02}e_j\bigr]
= s_0+(1-\lambda_{01})(1-\lambda_{02})e_1+\lambda_{02}^2e_j
\]
\[
\begin{cases}
(1-\lambda_{01})(1-\lambda_{02})=0 \\
(\lambda_{02})^2=1
\end{cases}
\]
\[
\boxed{\ \sigma_i s_j =
\begin{cases}
s_j & \text{si } j<i \\
s_{i+1} & \text{si } j=i \\
s_i & \text{si } j=i+1 \\
s_0+\lambda_i(e_i+e_{i+1})+e_j & \text{si } j\geq i+2
\end{cases}
\quad \text{on pose } e_0=0\ }
\]
\[
\boxed{\ \begin{aligned}
\sigma_0 e_j &=
\begin{cases}
-e_1 & j=1 \\
(\lambda_0-1)e_1+e_j & \text{si } j\geq 2
\end{cases} \\
\sigma_i e_j &=
\begin{cases}
e_j & \text{si } j<i \\
e_{i+1} & \text{si } j=i \\
e_i & \text{si } j=i+1 \\
\lambda_i(e_i+e_{i+1})+e_j & \text{autre (si } j\geq i+2)
\end{cases}
\quad (i\geq 1)
\end{aligned}\ }
\]
\note{les deux encadrés du bas donnent une forme déjà réduite à un paramètre $\lambda_i$ par réflexion, anticipant ce que les pp.~15--18 établissent.}

\page{12}
\[
\boxed{\ \sigma_i s_j = s_j \quad\text{si } j\leq i-1\ }
\qquad
\boxed{\ \begin{cases} \sigma_i s_i = s_{i+1} \\ \sigma_i s_{i+1} = s_i \end{cases}\ }
\]
\[
\sigma_i^2(s_j) = s_j \quad\text{si } j\leq i+1
\]
\[
\sigma_i^2 s_{i+2} = s_{i+2} \overset{?}{\Longrightarrow} \sigma_i^2 s_j = s_j \quad\text{si } j\geq i+3
\]
\[
\begin{aligned}
\sigma_i(\sigma_i s_j) &= \sigma_i(s_0+\lambda_{i1}e_1+\dots+\lambda_{i,i+1}e_{i+1}+\lambda_{i,i+2}e_j) \\
&= \underbrace{\sigma_i(s_0+\dots+\lambda_{i,i-1}e_{i-1})}_{s_0+\lambda_{i1}e_1+\dots+\lambda_{i,i-1}e_{i-1}} + (\lambda_{i,i}e_{i+1}+\lambda_{i,i+1}e_i) \\
&\quad + \lambda_{i,i+2}\underbrace{\sigma_i e_j}_{(\lambda_{i1}e_1+\dots+\lambda_{i,i+1}e_{i+1})+\lambda_{i,i+2}e_j} \qquad (i\geq 1) \\
&= s_0+(\lambda_{i1}+\lambda_{i,i+2}\lambda_{i1})e_1+(\lambda_{i2}+\lambda_{i,i+2}\lambda_{i2})e_2+\dots \\
&\quad +(\lambda_{i,i-1}+\lambda_{i,i+2}\lambda_{i,i-1})e_{i-1} \\
&\quad +(\lambda_{i,i+1}+\lambda_{i,i+2}\lambda_{i,i})e_i \\
&\quad +(\lambda_{i,i}+\lambda_{i,i+2}\lambda_{i,i+1})e_{i+1} \\
&\quad +(\lambda_{i,i+2})^2e_j
\end{aligned}
\]
\[
\sigma_i^2 s_j = s_j \iff
\begin{cases}
\lambda_{i1}(1+\lambda_{i,i+2}) = \lambda_{i2}(1+\lambda_{i,i+2}) = \dots = \lambda_{i,i-1}(1+\lambda_{i,i+2}) = 0 \\
\lambda_{i,i+1}+\lambda_{i,i+2}\lambda_{ii} = \lambda_{i,i}+\lambda_{i,i+2}\lambda_{i,i+1} = 0 \\
(\lambda_{i,i+2})^2 = 1
\end{cases}
\]
\struck{(\ill{})} pour $(n\geq j\geq i+2\geq i\geq 1)$.

\drawing{au bas de la page, esquisse d'une matrice par blocs : un bloc identité en haut à gauche (indice $i$ marqué au-dessus), des zéros au-dessous, suivie d'un second bloc à peine indiqué ; précédée d'un mot \ill{}.}

\page{13}
\[
\sigma_i s_j =
\begin{cases}
s_j & \text{si } j<i \\
s_{i+1} & \text{si } j=i \\
s_i & \text{si } j=i+1 \\
s_0+\lambda_{i,1}e_1+\dots+\lambda_{i,i+1}e_{i+1}+\lambda_{i,i+2}e_j & \text{si } i+2\leq j\leq n+1 \\
 & (\text{donc } 0\leq i\leq n-1)
\end{cases}
\]
$\begin{pmatrix} 0\leq i\leq n \\ 0\leq j\leq n+1 \end{pmatrix}$ \note{la parenthèse des bornes est en marge gauche.}
\[
\sigma_i e_j =
\begin{cases}
e_j & \text{si } j<i \\
e_{i+1} & \text{si } j=i \\
e_i & \text{si } j=i+1 \\
\lambda_{i,1}e_1+\dots+\lambda_{i,i+1}e_{i+1}+\lambda_{i,i+2}e_j & \text{si } i+2\leq j\leq n+1 \\
 & (\text{donc } 0\leq i\leq n-1)
\end{cases}
\qquad 1\leq i\leq n
\]
\[
\sigma_0 e_j =
\begin{cases}
-e_1 & \text{si } j=1 \\
(\lambda_{0,1}-1)e_1+\lambda_{0,2}e_j & \text{si } 2\leq j\leq n+1
\end{cases}
\]
\[
\sigma_0^2=\mathrm{id} \iff
\begin{cases}
(\lambda_{01}-1)(\lambda_{02}-1)=0 \\
\lambda_{02}^2=1
\end{cases}
\]
\[
\sigma_i^2=\mathrm{id} \ (1\leq i\leq n-1) \iff
\begin{cases}
\lambda_{i1}(1+\lambda_{i,i+2}) = \lambda_{i2}(1+\lambda_{i,i+2}) = \dots = \lambda_{i,i-1}(1+\lambda_{i,i+2}) = 0 \\
\lambda_{i,i+1}+\lambda_{i,i+2}\lambda_{ii} = \lambda_{i,i}+\lambda_{i,i+2}\lambda_{i,i+1} = 0 \\
(\lambda_{i,i+2})^2 = 1
\end{cases}
\]
\marginal{\emph{NB} ces équations \ill{} \ill{} si $i\geq 2$}

$\sigma_n^2=\mathrm{id}$ tjrs vérifié, $\sigma_n$ est une \emph{réflexion}.

Exprimons que $\sigma_0$ est une réflexion.
\note{en marge droite, une première matrice de $\sigma_0$, commencée (première ligne $-1,\ \lambda_{01}-1,\ \lambda_{01}-1\cdots$, première colonne de zéros), est biffée.}

\struck{$\sigma_0(x_1,\dots,x_{n+1}) = \sigma_0(s_0+x_1e_1+\dots+x_{n+1}e_{n+1})$}
\[
\sigma_0 =
\begin{pmatrix}
-1 & \lambda_{0,1}-1 & \lambda_{0,1}-1 & \cdots & \lambda_{01}-1 & 1 \\
0 & \lambda_{0,2} & 0 & & 0 & 0 \\
0 & 0 & \lambda_{02} & & 0 & 0 \\
\vdots & \vdots & \vdots & \ddots & \vdots & \vdots \\
0 & 0 & 0 & & \lambda_{02} & 0 \\
 & & & & 0 & 1
\end{pmatrix}
\]
\note{une accolade sous les colonnes porte « $n+1$ ».}
\[
\sigma_0(x_1,\dots,x_{n+1}) = \struck{\ill{}}\ \bigl(1-x_1+(\lambda_{0,1}-1)(x_2+\dots+x_{n+1}),\ \lambda_{0,2}x_2,\ \dots,\ \lambda_{0,2}x_{n+1}\bigr)
\]

\page{14}
Points fixes :
\[
H_0 :
\begin{cases}
(\lambda_{02}-1)x_i = 0 & \text{si } 2\leq i\leq n+1 \\
2x_1+(1-\lambda_{01})(x_2+\dots+x_{n+1}) = 1
\end{cases}
\]
\struck{\ill{}} Si $\lambda_{0,2}-1$ inv. ces équations s'écrivent \struck{$x_i=0$, $2\leq i\leq n+1$},
\[
\begin{cases}
x_i = 0 & 2\leq i\leq n+1 \\
2x_1\ \struck{\ill{}} = 1
\end{cases}
\]
donc c'est une sous-espace linéaire affine de dim.~$0$ \struck{\ill{}} \add{(en $\pm 2$)} \struck{\ill{}} si $2$ inv. \struck{\ill{}}, et c'est vide si $2$ nul dans $k$ \add{[\struck{\ill{}} \ill{} en car.~$2$, l'hyperplan : l'infini\ldots]} --- donc ce n'est pas un hyperplan. Donc il faut que $\lambda_{02}-1$ soit non inv. \add{i.e. nul} sur les fibres, donc nilpotent, i.e. $\lambda_{02}=1+\nu$, $\nu$ nilpotent. Alors les équations deviennent
\[
\begin{cases}
\nu x_i = 0 & 2\leq i\leq n+1 \\
2x_1+(1-\lambda_{01})(x_2+\dots+x_{n+1}) = 1
\end{cases}
\]
est un hyperplan sur chaque fibre, sauf \struck{en} en car.~$2$ si $\lambda_{0,1}=1$, où c'est vide (mais alors on trouve l'hyperplan : l'infini).
\marginal{introduction de la variable $2$ \ill{} \ill{} \ill{} $2x_1+(1-\lambda_{01})(x_2+\dots+x_{n+1})=2$}

Pour que ce soit un hyperplan, \struck{il faut \ill{}} \struck{\ill{}} \add{\ill{}} (où la car. est $\neq 2$ \struck{\ill{}} ou $\lambda_{01}-1$ inv.), comme \struck{l'hyperplan $H_0$} $H_0$ contenu dans $H_0'$ d'équation $2x_1+(1-\lambda_{01})(x_2+\dots+x_{n+1})=\uncertain{2}$, il faut qu'il y ait égalité, donc que les équations $\nu x_i=0$ \add{($2\leq i\leq n+1$)} soient conséquences de la précédente \add{(vu \ill{} de base\ldots)}, \add{d'évidence} ceci exige que $\nu=0$. \struck{Si $2$ \ill{} \ill{}. Mais $(1-\lambda_{01})$ \ill{}, \ill{} \ill{} \ill{} \ill{} \ill{} \ill{} \ill{} cas $n=1$ doit être traité à part\ldots\ Si $2$ est $(1-\lambda_{01})$}
\note{les trois dernières lignes sont biffées d'un trait ondulé.}

\page{15}
\struck{\ill{} inversibles, \ill{} \ill{} \ill{} \ill{} \ill{}, \ill{} \ill{} $(H_0')_t$ \ill{} hyperplan : l'infini\ldots}
\note{deux premières lignes biffées en diagonale.}

On exigera donc
\[
\lambda_{02} = 1
\]
(ce qui suffit à impliquer $\sigma_0^2=\mathrm{id}$) et on posera
\[
\boxed{\ \lambda_{01} = 1+\alpha_0\ }
\]
de sorte que $\alpha_0$ est l'invariant modulaire du polygone régulier induit par $\Pi$ dans $f_2$\ldots

\struck{Exprimons que $\sigma_i$ est une}
\[
\begin{cases}
\sigma_0 =
\begin{pmatrix}
-1 & \alpha_0 & \alpha_0 & \cdots & \alpha_0 & 1 \\
0 & 1 & 0 & \cdots & 0 & 0 \\
0 & 0 & 1 & \cdots & 0 & 0 \\
\cdots & & & & & \\
0 & 0 & 0 & \cdots & 1 & 0 \\
0 & 0 & 0 & \cdots & 0 & 1
\end{pmatrix} \\[2ex]
\sigma_0(x_1,\dots,x_{n+1}) = \bigl(1-x_1+\alpha_0(x_2+\dots+x_{n+1}),\ x_2,\ \dots,\ x_{n+1}\bigr) \\[1ex]
H_0 : 2x_1\ \uncertain{-}\ \alpha_0(x_2+\dots+x_{n+1}) = 1
\end{cases}
\]

Exprimons que $\sigma_i$ est une \emph{réflexion} ($1\leq i\leq n-1$).
\drawing{matrice par blocs de $\sigma_i$, $(n+2)\times(n+2)$ : bloc identité d'ordre $i-1$ en haut à gauche ; bloc $2\times 2$ $\begin{pmatrix}0&1\\1&0\end{pmatrix}$ aux places $i,i+1$ ; à droite, sur les lignes $1$ à $i+1$, les coefficients $\lambda_{i1}\ \lambda_{i1}\cdots\lambda_{i1}$, \ldots, $\lambda_{ii}\cdots\lambda_{ii}$, $\lambda_{i,i+1}\cdots\lambda_{i,i+1}$ ; sur les $n-i$ dernières lignes, la diagonale $\lambda_{i,i+2}$ ; dernière colonne de translation nulle, $1$ en bas à droite. Les accolades marquent « $i-1$ », « $2$ », « $n-i$ », « $i+1$ ».}
\[
\begin{aligned}
\sigma_i(x_1,\dots,x_{n+1}) = \bigl(&x_1+\lambda_{i1}(x_{i+2}+\dots+x_{n+1}),\ x_2+\lambda_{i2}(x_{i+2}+\dots+x_{n+1}),\ \dots, \\
&x_{i-1}+\lambda_{i,i-1}(x_{i+2}+\dots+x_{n+1}),\ x_{i+1}+\lambda_{i,i}(x_{i+2}+\dots+x_{n+1}), \\
&x_i+\lambda_{i,i+1}(x_{i+2}+\dots+x_{n+1}),\ \lambda_{i,i+2}x_{i+2},\ \dots,\ \lambda_{i,i+2}x_{n+1}\bigr)
\end{aligned}
\]

\page{16}
\note{le haut de la page, jusqu'à « $\lambda_{i,i+2}=1$ », est biffé de longues diagonales ; il est transcrit.}
\[
\struck{\sigma_i(x_1,\dots,x_{n+1}) = (x_1,\dots,x_{i-1},x_{i+1},x_i,\lambda_{i,i+2}x_{i+2},\lambda_{i,i+2}x_{i+3},\dots,\lambda_{i,i+2}x_{n+1})}
\]
\[
\struck{H_i : \begin{cases} x_i = x_{i+1} \\ (\lambda_{i,i+2}-1)x_{i+2} = \dots = (\lambda_{i,i+2}-1)x_{n+1} = 0 \end{cases}}
\]
\struck{Si $H_i'$ est l'hyperplan d'équation $x_i=x_{i+1}$, on a $H_i\subset H_i'$, donc $H_i$ n'est un hyperplan que si $H_i=H_i'$, i.e. les équations $(\lambda_{i,i+2}-1)x_{i+2}=\dots=(\lambda_{i,i+2}-1)x_{n+1}=0$ sont conséquences de $x_i=x_{i+1}$, ce qui \ill{} \ill{} que si $\lambda_{i,i+2}=1$.}
\[
(H_i)\quad
\begin{cases}
\lambda_{i1}(x_{i+2}+\dots+x_{n+1}) = \lambda_{i2}(x_{i+2}+\dots+x_{n+1}) = \dots \\
\qquad = \lambda_{i,i-1}(x_{i+2}+\dots+x_{n+1}) = 0 \\
\lambda_{i,i}(x_{i+2}+\dots+x_{n+1}) = x_i-x_{i+1} \\
\lambda_{i,i+1}(x_{i+2}+\dots+x_{n+1}) = x_{i+1}-x_i \\
(\lambda_{i,i+2}-1)x_{i+2} = \dots = (\lambda_{i,i+2}-1)x_{n+1} = 0
\end{cases}
\]
(en plus, les conditions sur les $\lambda_{i,j}$, $1\leq j\leq i+2$, \struck{\ill{}} exprimant $\sigma_i^2=\mathrm{id}$\ldots)

\uncertain{Si par exemple} : si on avait $\lambda_{i,i+2}\neq 1$, comme $\lambda_{i,i+2}^2=1$, on aurait $\lambda_{i,i+2}=-1\neq 1$ (car $\neq 2$), et les équations de $H_i$ deviennent
\[
\begin{cases}
x_{i+2} = \dots = x_{n+1} = 0 \\
\struck{\ill{} = \lambda_{i,i+1}\Sigma_i(x) = 0 \quad \Sigma_i(x)=x_{i+2}+\dots+x_{n+1}} \\
\struck{\ill{}\ x_i = x_{i+1}} \\
\struck{\ill{}}
\end{cases}
\]
ce qui fait $[n+1-(i+1)]+1 = n-i+1$ équations \struck{si $i\leq n$, les premières équations \ill{} \ill{}} indépendantes et comme $i\leq n-1$, on trouve une variété linéaire de codim.~$\geq 2$. Donc on doit avoir $\lambda_{i,i+2}=1$, les équations de $H_i$ deviennent

\page{17}
\[
\begin{cases}
\lambda_{i1}\Sigma_i(x) = \dots = \lambda_{i,i-1}\Sigma_i(x) = 0 \\
\lambda_{i,i}\Sigma_i(x) = x_i-x_{i+1} \\
\lambda_{i,i+1}\Sigma_i(x) = x_{i+1}-x_i
\end{cases}
\qquad \text{où } \Sigma_i(x) = x_{i+2}+\dots+x_{n+1}
\]
Si dans les $\lambda_{ij}$, $1\leq j\leq i-1$, un est $\neq$ nul, alors ces équations se réduisent à :
\[
\begin{cases}
\Sigma_i(x) = 0 \\
x_i-x_{i+1} = 0
\end{cases}
\]
ce sont des équations linéaires indépendantes, donc $H_i$ serait de codim.~$2$, donc pas un hyperplan. On veut donc
\[
\lambda_{i1} = \dots = \lambda_{i,i-1} = 0
\]
et les équations se réduisent alors à :
\[
\lambda_{ii}\Sigma_i(x) = -\lambda_{i,i+1}\Sigma_i(x) = x_i-x_{i+1} .
\]
\struck{Si on avait $\lambda_{ii}=\lambda_{i,i+1}=0$, ces équations}
Les relations exprimant $\sigma_i^2=\mathrm{id}$ se réduisent ici à : $\lambda_{i,i}+\lambda_{i,i+1}=0$, de sorte qu'il reste un seul paramètre $\lambda_{i,i+1}=1+\alpha_i$, et on trouve
\drawing{matrice par blocs de $\sigma_i$ : identité d'ordre $i-1$ en haut à gauche ; bloc $\begin{pmatrix}0&1\\1&0\end{pmatrix}$ aux places $i,i+1$, suivi sur ces deux lignes de $-(1+\alpha_i)\ \cdots\ -(1+\alpha_i)$ et $(1+\alpha_i)\ \cdots\ (1+\alpha_i)$ ; identité sur le reste de la diagonale ; dernière colonne $(0,\dots,0,1)$. Plusieurs entrées retouchées.}

\page{18}
\[
\begin{cases}
\begin{aligned}
\sigma_i(x_1,\dots,x_{n+1}) = \bigl(&x_1,\dots,x_{i-1},\ x_{i+1}-(1+\alpha_i)(x_{i+2}+\dots+x_{n+1}), \\
&x_i+(1+\alpha_i)(x_{i+2}+\dots+x_{n+1}), \\
&x_{i+2},\dots,x_{n+1}\bigr)
\end{aligned} \\[2ex]
H_i :\ x_i-x_{i+1}+(1+\alpha_i)(x_{i+2}+\dots+x_{n+2}) = 0
\end{cases}
\]
\note{« $x_{n+2}$ » est écrit ainsi, pour $x_{n+1}$.}

Ces résultats marchent aussi sur un anneau réduit. Pour un anneau quelconque [il y a lieu de faire \uncertain{poser}]
\[
\lambda_{i,i+2} = 1+\nu \qquad \nu \text{ nilpotent}
\]
\struck{et les relations} en car. rés. $\neq 2$ les relations qui expriment $\sigma_i^2=\mathrm{id}$ deviennent $\nu$
\[
\begin{cases}
\nu = 0 \\
\lambda_{i,1} = \dots = \lambda_{i,i-1} = 0 \\
\lambda_{i,i+1}+\lambda_{i,i+1}\ \struck{\ill{}} = \lambda_{i,i}+\lambda_{i,i+1}\ \struck{\ill{}} = 0 \\
\struck{\nu^2=\ill{}}
\end{cases}
\]
\struck{donc $\lambda_{i,i}+\lambda_{i,i+1} = -\nu\lambda_{ii} = -\nu(\lambda_{i,i+1})$}

Comme devant. Il reste à examiner ce qui se passe en car. rés.~$2$ (ou un anneau quelconque, si on veut\ldots), où on trouve :
\[
\begin{cases}
\underbrace{(2+\nu)}_{\text{nilpot.}}\lambda_{i,1} = (2+\nu)\lambda_{i,2} = \dots = (2+\nu)\lambda_{i,i-1} = 0 \\
\lambda_{ii}+\lambda_{i,i+1} = -\nu\lambda_{ii} = -\nu\lambda_{i,i+1} \\
\struck{\ill{}}\ \nu^2+2\nu = 0 \quad \text{i.e.}\quad \nu(2+\nu) = 0
\end{cases}
\]
\[
\Longrightarrow\quad \lambda_{i,i+1} = -(1+\nu)\lambda_{ii} \qquad \struck{-\nu\lambda_{ii}=}\quad \nu\underbrace{\bigl[\lambda_{i,i+1}-\lambda_{ii}\bigr]}_{-(2+\nu)\lambda_{ii}} = 0
\]
\note{la troisième ligne du premier système, en car. rés. $\neq 2$, se lit $\lambda_{i,i+1}+\lambda_{i,i+1}$ ; on attendrait $\lambda_{i,i+1}+\lambda_{i,i}$.}

\page{19}
\[
\text{II}\qquad
\left.\begin{cases}
(2+\nu)\lambda_{i,1} = \dots = (2+\nu)\lambda_{i,i-1} = 0 \\
\lambda_{i,i+1} = -(1+\nu)\lambda_{i,i} \\
\nu(2+\nu) = 0
\end{cases}\right]
\]
Reprenons les équations de $H_i$. $H_i$ est contenu dans l'hyperplan $H_i'$ d'équation
\[
\lambda_{ii}\Sigma_i(x) = x_i-x_{i+1} .
\]
\struck{Si on avait} \uncertain{Si} $H_i$ était un hyperplan, on aurait $H_i=H_i'$, i.e. les autres équations de $H_i$ seraient conséquences de la précédente. \uncertain{Ce n'est possible} manifestement que si $\lambda_{i,1}=\dots=\lambda_{i,i-1}=0$, $\lambda_{i,i}+\lambda_{i,i+1}=0$, $\lambda_{i,i+1}=1$, OK\ldots

\emph{NB} $\sigma_0$ ne change que la première coordonnée, $\sigma_i$ ($1\leq i\leq n-1$) que la $i$-ième et $(i+1)$-ième. Ceci \struck{\ill{}} \add{implique} que $\sigma_i$ et $\sigma_j$ commutent si $j\geq i+2$.
\[
\begin{aligned}
\sigma_0\sigma_1(x_1,x_2,x_3,\dots,x_{n+1}) &= \sigma_0\bigl(x_2-(1+\alpha_1)\Sigma_1,\ x_1+(1+\alpha_1)\Sigma_1,\ x_3,\dots\bigr) \\
&= \bigl(1-x_2+(1+\alpha_1)\Sigma_1+\alpha_0(x_1+(1+\alpha_1)\Sigma_1+\Sigma_1), \\
&\qquad x_1+(1+\alpha_1)\Sigma_1,\ x_3,\dots,x_{n+1}\bigr) \\
&= \bigl(1+\alpha_0x_1-x_2+(1+\alpha_1+2\alpha_0+\alpha_0\alpha_1)\Sigma_1,\ x_1+(1+\alpha_1)\Sigma_1, \\
&\qquad x_3,\dots,x_{n+1}\bigr)
\end{aligned}
\]
\note{dans la dernière ligne, un premier coefficient de $\Sigma_1$ est biffé et remplacé au-dessous par $(1+\alpha_1+2\alpha_0+\alpha_0\alpha_1)$.}
\[
(\sigma_0\sigma_1)^p(x_1,x_2,\dots,x_{n+1}) = \bigl(F_p(\alpha_0,\alpha_1;x_1,\dots,x_{n+1}),\ G_p(\alpha_0,\alpha_1;x_1,\dots,x_{n+1}),\ x_3,\dots,x_{n+1}\bigr)
\]
Donc $(\sigma_0\sigma_1)^p = \mathrm{id} \iff$
\note{la phrase s'arrête là ; la page suivante reprend sous une autre forme (restrictions aux quotients du drapeau) et le calcul de $(\sigma_0\sigma_1)^p$ n'est pas mené à terme dans ce lot.}

\page{20}
\[
(i\geq 1)\quad
\begin{cases}
\sigma_i|f_{i-1} = \mathrm{id} \\
\sigma_i|(f_{i+1}/f_{i-1}) = \text{échange des coordonnées } x_i,x_{i+1} \\
\sigma_i|E/f_{i+1} = \mathrm{id}
\end{cases}
\]
\[
\begin{cases}
\sigma_0|f_1 = \text{échange de } s_0,s_1 \\
\sigma_0|E/f_1 = \mathrm{id}
\end{cases}
\]
\[
\struck{\ill{}}\ (i\geq 1)\quad
\begin{cases}
\sigma_i\sigma_{i+1}|f_{i-1} = \mathrm{id} \\
\sigma_i\sigma_{i+1}|f_{i+2}/f_{i-1} \overset{?}{=} \\
\sigma_i\sigma_{i+1}|E/f_{i+2} = \mathrm{id}
\end{cases}
\]
\[
\begin{cases}
\sigma_0\sigma_1|f_2 \overset{?}{=} \\
\sigma_0\sigma_1|E/f_2 = \mathrm{id}
\end{cases}
\]
\note{les deux « $\overset{?}{=}$ » restent sans membre de droite ; le reste de la page est blanc. La liasse annoncée « 23~p » sur la couverture se poursuit au-delà de ce lot.}

\end{document}
